\documentclass[11pt]{article}

\usepackage[margin=1in]{geometry}
\usepackage{hyperref}
\usepackage{enumitem}
\usepackage{booktabs}
\usepackage{amsmath}
\usepackage{listings}
\usepackage{amssymb}
% Optional: Configure listings style
\usepackage{tikz}

\usepackage{amsmath,amssymb,amsfonts}
\usepackage{algorithmic}
\usepackage{graphicx}
\usepackage{textcomp}
\usepackage{xcolor}
\usepackage{listings}
\usepackage{xcolor}

% Optional: Configure listings style
\lstset{
    basicstyle=\ttfamily\small,
    keywordstyle=\color{blue},
    commentstyle=\color{gray},
    stringstyle=\color{red},
    showstringspaces=false,
    breaklines=true,
    frame=single,
    numbers=left,
    numberstyle=\tiny\color{gray},
    captionpos=b
}

\title{QLego Implementation}
\author{Abhishek Shringi}
\date{January 2026}

\begin{document}
\maketitle
The design principles behind QLego are simple viz. modular by nature and extensible by design. In order to benchmark the passes across the compiler frameworks, such a design philosophy comes in handy.
\section{Overall Framework}
The framework is primarily divided into 2 components, namely Core and Plugin.
\section{Core}
The primary responsibility of Core is to orchestrate the flow of various passes in the pipeline. It also handles the environment swtiching across the different plugin's pass
\subsection{QPass}
QPass is an abstract, with the purpose of handling the execution of pass from a particular plugin written with in a specific framework. QPass also provides an abstract method "run" which particular user/plugin can configure to make their passes run in a particular manner based on whatever the particular framework allows.
\subsection{QPipeline}
QPipeline is the actual compiler here, which brings together various QPass( could be from different Plugins ) and eventually executes them in a sequential manner. Further, it also has capability to aggregate adjacent QPass, if they have no dependency conflicts, thereby saving the unnecessary overhead cost while switching the environment. One thing to note here is that while QPipeline allows us to do compilation in form of standard stages, it has the flexibility to experiment unconventional compilation flows like merging routing and mapping.
\subsection{QPassContext}
QPassContext is basically abstract for the context object which is shared across various QPass in the QPipeline. It consists of following components:
\begin{itemize}
    \item \textbf{qasm:} OpenQASM string which contains the current state of circuit. Currently, OpenQASM2 has been chosen as the exchange format of the circuit across passes as it is widely supported by multiple frameworks.
    \item \textbf{hardware:} This contains the backend information in form of a json string. More details can be found in the subsequent section
    \item \textbf{metadata:} This can contain arbitrary data related to analysis with in various QPass
\end{itemize}
 
\subsection{QBackend}
Every framework has different support for its backend, be it a real machine or noisy simulator, for instance Qiskit has got FakeBackendV2, AerSimulator or BackendV2. Therefore, there is a need to abstract out the backend used in various frameworks as this information could be required by passes like mapping and routing.
\section{Plugin}
Plugin though a common term, here has been used to highlight the extensibility of QLego. A plugin corresponding to a framework consists of following components
\subsection{Adapter}
This contains all the required things to port and run circuit pass in a particular framework.
\subsection{Environment}
This contains the SDK itself along with required dependencies in an isolated enviroment. The isolation itself is a abstract concept and could be at various levels. Here, we make provide isolations at two levels. Since, most of the Quantum SDKs are pythonic in nature, we provide "venv" isolation for those SDKs. Furthermore, some of the SDKs have native C++ dependencies, which could be handled by "conda" environment isolation.
\section{Protocol}
Here, we demonstrate the flow of QLego by an example of an Bell Pair generation circuit. The input is given in form of OpenQasm 2. The context is generated by passing the input circuit and the backend in whatever format we like, QBackend will take care of conversion to JSON and passing it to constituent QPass in QPipeline. The QPipeline is instantiated with the required passes, each QiskitPass and TKetPass are the ?instances? of abstract QPass. When the pipeline is run, it goes sequentially across the pipeline over each pass. Each pass( QiskitPass and TketPass ) have venv path, so that pipeline knows where to execute that pass. Each pass is run as a separate process in its own environment. And the output finally is stored in the QPassContext object and final circuit could be retreived by calling the member "qasm". The "metadata" member would contain information like runtime of each pass and circuit metrics if you would have added evaluation pass as well.
\begin{figure}[!t]
\begin{lstlisting}[
    caption={Example usage of the QLego framework to construct a hybrid compilation pipeline combining Qiskit and TKet passes.}, 
    label={lst:usage}, 
    language=Python,
    basicstyle=\ttfamily\small,
    keywordstyle=\color{blue},
    commentstyle=\color{gray},
    stringstyle=\color{red},
    showstringspaces=false,
    breaklines=true,
    frame=single,
    numbers=none,
    captionpos=b,
    xleftmargin=0pt,
    xrightmargin=0pt
]
from qlego.qpass import QPipeline, QPassContext
from qlego_qiskit.adapter.passes import QiskitPass
from qlego_tket.adapter.passes import TKetPass
from qiskit.transpiler.passes import Optimize1qGates, CXCancellation

pipeline = QPipeline([
    QiskitPass([
        Optimize1qGates(),
        CXCancellation()
    ]),
    
    TKetPass([
        "PlacementPass",
        "RoutingPass"
    ])
])

qasm_input = """
OPENQASM 2.0;
include "qelib1.inc";
qreg q[2];
h q[0];
cx q[0], q[1];
"""

ctx = QPassContext(qasm=qasm_input)
result = pipeline.run(qasm_input, ctx)

print("Compiled circuit:")
print(result.qasm)
\end{lstlisting}
\end{figure}


\end{document}